% GENERATED FILE. Edit canonical lessons, then run npm run generate:books.
% canonical-source-hash: fnv1a64:6e0658032847b6a9
% canonical-lessons: LA-C464-statera, LA-C464-trutina, LA-C464-lanx, LA-C464-appendo, LA-C464-penso

\chapter{Steelyard, Pair of scales, Scale pan, To weigh out, To weigh carefully}
\label{ch:la-steelyard-pair-of-scales-scale-pan-to-weigh-out-to-weigh-carefully}
\hlchaptermodality{464}

\begin{chapteropening}
\textbf{By the end of this chapter:} \emph{I can talk about weighing in the market: the steelyard, a pair of scales, the scale pan, and weighing things out with care.}

Complete the last lesson of chapter 464: I can talk about weighing in the market: the steelyard, a pair of scales, the scale pan, and weighing things out with care.
\end{chapteropening}

\section[statēra, statērae]{statēra, statērae — a steelyard, a balance with one arm}

\label{lesson:LA-C464-statera}



\begin{quote}
Before the new one: say the Latin for a wicker basket, then the Latin for a bread basket.
\end{quote}

\subsection*{You'll want to know: statēra, statērae}
\textbf{statēra, statērae} — \textquotedblleft{}a steelyard, a balance with one arm\textquotedblright{}.

\begin{quote}
Lanius carnem in statērā pōnit.
\end{quote}

\begin{quote}
\emph{The butcher puts the meat on the steelyard.}
\end{quote}

First declension, feminine: \emph{statēra, statērae f.} A beam hung from a hook, with a weight slid along one arm.

\begin{cousinweb}
\emph{statēra} is usually traced to Greek \emph{statēr}, a weight and a coin, built on \emph{histēmi} 'to make stand, to weigh' (EtymDB). Lewis \& Short compare Greek \emph{stathmē} 'rule, balance' instead, so the exact source is not settled.
\end{cousinweb}

\begin{grammarlens}[title={What you've built}]
One more word for telling, explaining and describing at B1.
\end{grammarlens}

\subsection*{Your turn}
\begin{itemize}
\raggedright
  \item \emph{Say it:} \emph{statēra, statērae}
  \item \emph{Say it:} \emph{statēra, statērae}, once more
  \item \emph{Say it:} a sentence of your own with \emph{statēra, statērae}
  \item \emph{Recall:} say the Latin for a wicker basket, then the Latin for a bread basket, then say \emph{statēra, statērae} again
\end{itemize}

\begin{tcolorbox}[breakable,colback=teal!4,colframe=teal!35!black,title={Before you move on}]
What does statēra, statērae mean? (\textquotedblleft{}A steelyard, a balance with one arm\textquotedblright{}.) Say it once more.
\end{tcolorbox}
\section[trutina, trutinae]{trutina, trutinae — a pair of scales, a balance}

\label{lesson:LA-C464-trutina}



\begin{quote}
Before the new one: say the Latin for a bread basket, then the Latin for a steelyard.
\end{quote}

\subsection*{You'll want to know: trutina, trutinae}
\textbf{trutina, trutinae} — \textquotedblleft{}a pair of scales, a balance\textquotedblright{}.

\begin{quote}
Mercātor aurum in trutinā dīligenter pōnit.
\end{quote}

\begin{quote}
\emph{The merchant carefully puts the gold on the scales.}
\end{quote}

First declension, feminine: \emph{trutina, trutinae f.} Also figurative: weighing a thing in the mind.

\begin{cousinweb}
\emph{trutina} was borrowed from Greek \emph{trytanē} 'balance, tongue of a balance' (Lewis \& Short). EtymDB lists it without a parent.
\end{cousinweb}

\begin{grammarlens}[title={What you've built}]
One more word for telling, explaining and describing at B1.
\end{grammarlens}

\subsection*{Your turn}
\begin{itemize}
\raggedright
  \item \emph{Say it:} \emph{trutina, trutinae}
  \item \emph{Say it:} \emph{trutina, trutinae}, once more
  \item \emph{Say it:} a sentence of your own with \emph{trutina, trutinae}
  \item \emph{Recall:} say the Latin for a bread basket, then the Latin for a steelyard, then say \emph{trutina, trutinae} again
\end{itemize}

\begin{tcolorbox}[breakable,colback=teal!4,colframe=teal!35!black,title={Before you move on}]
What does trutina, trutinae mean? (\textquotedblleft{}A pair of scales, a balance\textquotedblright{}.) Say it once more.
\end{tcolorbox}
\section[lanx, lancis]{lanx, lancis — a scale pan; a dish, a platter}

\label{lesson:LA-C464-lanx}



\begin{quote}
Before the new one: say the Latin for a steelyard, then the Latin for a pair of scales.
\end{quote}

\subsection*{You'll want to know: lanx, lancis}
\textbf{lanx, lancis} — \textquotedblleft{}a scale pan; a dish, a platter\textquotedblright{}.

\begin{quote}
In ūnā lance merx est, in alterā pondus.
\end{quote}

\begin{quote}
\emph{In one scale pan there are the goods, in the other the weight.}
\end{quote}

Third declension, feminine: \emph{lanx, lancis f.} English \emph{balance} comes, through French, from \emph{bilanx} 'having two pans'.

\begin{cousinweb}
\emph{lanx} is inherited; Lewis \& Short and EtymDB compare Greek \emph{lekos} 'dish, pan' as a cousin. English \emph{balance} comes, through French, from Late Latin \emph{bilanx} 'with two scale pans'.
\end{cousinweb}

\begin{grammarlens}[title={What you've built}]
One more word for telling, explaining and describing at B1.
\end{grammarlens}

\subsection*{Your turn}
\begin{itemize}
\raggedright
  \item \emph{Say it:} \emph{lanx, lancis}
  \item \emph{Say it:} \emph{lanx, lancis}, once more
  \item \emph{Say it:} a sentence of your own with \emph{lanx, lancis}
  \item \emph{Recall:} say the Latin for a steelyard, then the Latin for a pair of scales, then say \emph{lanx, lancis} again
\end{itemize}

\begin{tcolorbox}[breakable,colback=teal!4,colframe=teal!35!black,title={Before you move on}]
What does lanx, lancis mean? (\textquotedblleft{}A scale pan; a dish, a platter\textquotedblright{}.) Say it once more.
\end{tcolorbox}
\section[appendō, appendere]{appendō, appendere — to weigh out; to hang on}

\label{lesson:LA-C464-appendo}



\begin{quote}
Before the new one: say the Latin for a pair of scales, then the Latin for a scale pan.
\end{quote}

\subsection*{You'll want to know: appendō, appendere}
\textbf{appendō, appendere} — \textquotedblleft{}to weigh out; to hang on\textquotedblright{}.

\begin{quote}
Pistor pānem appendit et pretium dīcit.
\end{quote}

\begin{quote}
\emph{The baker weighs out the bread and says the price.}
\end{quote}

\emph{appendō, appendere, appendī, appēnsum}. Money was once weighed, so \emph{appendere} can also mean to pay. English \emph{append} and \emph{appendix} come from it.

\begin{cousinweb}
\emph{appendō} is built from \emph{ad-} 'to' and \emph{pendō} 'to hang, to weigh' (EtymDB; Lewis \& Short). English \emph{append} and \emph{appendix} come from it.
\end{cousinweb}

\begin{grammarlens}[title={What you've built}]
One more word for telling, explaining and describing at B1.
\end{grammarlens}

\subsection*{Your turn}
\begin{itemize}
\raggedright
  \item \emph{Say it:} \emph{appendō, appendere}
  \item \emph{Say it:} \emph{appendō, appendere}, once more
  \item \emph{Say it:} a sentence of your own with \emph{appendō, appendere}
  \item \emph{Recall:} say the Latin for a pair of scales, then the Latin for a scale pan, then say \emph{appendō, appendere} again
\end{itemize}

\begin{tcolorbox}[breakable,colback=teal!4,colframe=teal!35!black,title={Before you move on}]
What does appendō, appendere mean? (\textquotedblleft{}To weigh out; to hang on\textquotedblright{}.) Say it once more.
\end{tcolorbox}
\section[pēnsō, pēnsāre]{pēnsō, pēnsāre — to weigh carefully, to weigh out; to ponder}

\label{lesson:LA-C464-penso}



\begin{quote}
Before the new one: say the Latin for a scale pan, then the Latin for to weigh out.
\end{quote}

\subsection*{You'll want to know: pēnsō, pēnsāre}
\textbf{pēnsō, pēnsāre} — \textquotedblleft{}to weigh carefully, to weigh out; to ponder\textquotedblright{}.

\begin{quote}
Mercātor prūdēns mercem pēnsat.
\end{quote}

\begin{quote}
\emph{A careful merchant weighs the goods.}
\end{quote}

\emph{pēnsō, pēnsāre, pēnsāvī, pēnsātum}, a frequentative of \emph{pendō}. English \emph{pensive} comes from it through French \emph{penser} 'to think'.

\begin{cousinweb}
\emph{pēnsō} is a frequentative of \emph{pendō} 'to weigh', made from its participle \emph{pēnsus} (Lewis \& Short). English \emph{pensive} comes from it through French \emph{penser} 'to think, to weigh in the mind'.
\end{cousinweb}

\begin{grammarlens}[title={What you've built}]
One more word for telling, explaining and describing at B1.
\end{grammarlens}

\subsection*{Your turn}
\begin{itemize}
\raggedright
  \item \emph{Say it:} \emph{pēnsō, pēnsāre}
  \item \emph{Say it:} \emph{pēnsō, pēnsāre}, once more
  \item \emph{Say it:} a sentence of your own with \emph{pēnsō, pēnsāre}
  \item \emph{Recall:} say the Latin for a scale pan, then the Latin for to weigh out, then say \emph{pēnsō, pēnsāre} again
\end{itemize}

\begin{tcolorbox}[breakable,colback=teal!4,colframe=teal!35!black,title={Before you move on}]
What does pēnsō, pēnsāre mean? (\textquotedblleft{}To weigh carefully, to weigh out; to ponder\textquotedblright{}.) Say it once more.
\end{tcolorbox}
